% Chapter 4.
\chapter{Conclusions}
\label{ch:conclusions}

\section{What operational data can tell us}
\label{sec:summary}

This thesis set out to learn the state of health of grid-scale batteries
from the data their operators already collect. Its central finding is
that this is possible, provided the learning is anchored to a physical
ageing law. Purely data-driven models extrapolate badly to ages and
conditions not seen in training \parencite{severson2019data}; purely physical models carry systematic
error from the laboratory into the field, however carefully their parameters are tracked \parencite{plett2004ekf}. Together, the physics carries
the trend and the data correct its shape.

Two results are worth carrying beyond the specific estimator. The first
is that, under frequency-response duty, batteries age with the calendar
more than with use (\cref{kr:calendar}), so the warranty unit of
equivalent full cycles should be supplemented by a measure of time spent
at high state of charge and temperature. The second is that an estimate
with honest uncertainty has direct economic value: the field trial
recovered usable energy without detectable cost in ageing.

\section{Limitations and further work}
\label{sec:limits}

The conclusions rest on eleven sites and two chemistries over three
years, a short span against a design life of fifteen. In particular:

\begin{itemize}
  \item the square-root calendar law may not hold once the
    interphase growth gives way to other mechanisms late in life, and
    the estimator has not been tested beyond 12\,\% capacity loss;
  \item the partial charge features require segments of at least twenty
    minutes, which sites dedicated to fast frequency response may never
    provide;
  \item the field trial ran on a single site, in a mild year, and its
    economic result depends on market prices that have since fallen.
\end{itemize}

Each of these points is a question that time and more sites will answer.
The data set and the estimator are released with this thesis so that
others can ask them.
